% Related work section fragment — natbib \citep/\citet; no preamble.
% Every citation key resolves to an entry in references.bib.
\section{Related work}
\label{sec:related}

Cell-state engineering poses an inverse-design question: given a starting state and a
desired target state, which interventions move the cell there, and is the move even
possible? This is the question behind directed differentiation, transdifferentiation, and
T-cell engineering, and it is the question of pharmaceutical target selection. Getting it
wrong is the dominant cost of drug development: only about one in ten clinical programmes
reaches approval, and failure is paid for late \citep{Minikel2024,WongSiahLo2019}. Attrition concentrates
at Phase~II, and roughly half of Phase~II/III failures are for lack of efficacy---the
wrong target or mechanism rather than safety \citep{Harrison2016}. The two interventions
with published evidence of improving those odds both point the same way: human genetic
support raises the odds of approval about 2.6-fold \citep{Minikel2024,Nelson2015}, and
\emph{measured} rather than inferred target effects are the direct antidote to the
11--25\% preclinical reproducibility documented by the Amgen and Bayer audits
\citep{BegleyEllis2012,Prinz2011}. A method for this problem should therefore (i)~work from
measured effects in the relevant cell type, (ii)~answer achievability rather than merely
rank candidates, and (iii)~say so honestly when a goal is not achievable.

We surveyed 91 prior methods spanning 2011--2026 (92 including this work) across
control-theoretic network control,
gene-regulatory-network (GRN) inference with in-silico perturbation, deep
perturbation-response prediction, single-cell foundation models, and the
combinatorial-screen, optimal-transport, virtual-cell, and benchmark literature that
borders them. Positioning them on two axes---\emph{what question is answered} (forward
prediction versus feasibility) and \emph{what the answer is grounded in} (an inferred or
hand-built model versus directly measured effects)---exposes a gap that is exact rather
than rhetorical. The methods grounded in measured effects answer only forward questions,
and the methods that reason about feasibility do so only inside a model.

\paragraph{Model-theoretic feasibility.}
The most mature theory for the intervention question is network control. Structural and
target controllability cast steering an assumed linear system as a matching problem over
network topology \citep{Liu2011,Gao2014}, and determining-node / control-kernel methods
identify driver sets from structure \citep{Mochizuki2013,Kim2013}. Boolean attractor
control computes node interventions that switch a logical network between attractors
\citep{Zanudo2015,Zanudo2017,Joo2018,Marazzi2022}, a line that remains active
\citep{Cifuentes-Fontanals2025,Metzcar2025,Murrugarra2025,Trinh2025,Vecchio2026}. A
data-fitted strand fits a dynamical model to omics and then applies control theory
\citep{Ronquist2017,Rukhlenko2022,Wytock2024}, and network-influence predictors rank
transcription-factor sets for a conversion \citep{Rackham2016,Cahan2014}. These are the
methods that carry a controllability or control-strategy guarantee---but the guarantee is
a property of a hand-built, fitted, or inferred model, never of measured effects, and the
control set is only as trustworthy as the network put in. The same dependency reappears
whenever feasibility-style reasoning is imported onto a reconstructed single-cell GRN
\citep{Wang2023CEFCON} or run in a learned latent space \citep{Han2025}. The set of prior
methods that emit any verdict- or certificate-like output is therefore not one family: it
spans structural and Boolean network control, control-from-data, an inferred single-cell
GRN, and one causally inspired perturbation predictor. What unites them is the assumed
model, not the family.

\paragraph{Measured-grounded forward prediction.}
The single-cell community answers the same biological question by running an inferred
network forward. CellOracle infers a GRN and simulates a transcription-factor knockout
into predicted identity transitions \citep{Kamimoto2023}; related tools infer regulons,
enhancer networks, or dynamic GRNs \citep{Aibar2017,Xu2021,BravoGonzalezBlas2023,%
Wang2023Dictys,Osorio2022}, with BEELINE quantifying how method-dependent GRN inference
is \citep{Pratapa2020}. Deep perturbation-response models are genuinely
measured-grounded---they learn from real Perturb-seq data---but are resolutely forward:
GEARS predicts outcomes of unseen multi-gene perturbations \citep{Roohani2023}, and
compositional and disentanglement models predict responses to given perturbations
\citep{Lotfollahi2019,Lotfollahi2023,Hetzel2022,Piran2024}. Single-cell foundation models
\citep{Theodoris2023,Cui2024,Hao2024} and the 2025--2026 virtual-cell surge---perturbation
transformers, atlas-scale models and challenges, and giga-scale perturbation atlases
\citep{Adduri2025,Roohani2025,Zhang2025Tahoe,Klein2025b}---extend the same forward map.
Optimal-transport methods describe how cells \emph{do} move rather than which interventions
\emph{would} move them \citep{Schiebinger2019,Lange2022}, and experimental-design methods
choose the next screen \citep{Huang2024}. Every one maps perturbation to predicted
response: a predictor always returns \emph{some} answer and structurally cannot certify
that a target is unreachable.

\paragraph{Why a feasibility verdict, not another predictor.}
A cluster of 2025--2026 benchmarks reports that current perturbation-prediction models,
foundation models included, often fail to beat simple mean or linear baselines on
genuinely unseen perturbations \citep{AhlmannEltze2025,Wong2025,Csendes2025,%
Kedzierska2025,Wu2025,Vinas2025,Kernfeld2025,Wei2025,Luecken2025,Atti2025,Mao2026,%
Vollenweider2026}. This is the field's own critique, and it sharpens the case for a
different output: a feasibility verdict computed directly from measured effects does not
extrapolate to unseen perturbations at all, and it comes with a checkable certificate and a
shuffled-target null. The claim is deliberately weaker, and therefore more defensible, than
predicting an unseen response.

\paragraph{Our position.}
Across all 91 prior methods (ours excluded), 14 operate directly on measured effects---and
all 14 only predict or rank; \emph{none} returns a feasibility verdict, emits a
data-grounded certificate, or proves infeasibility, and the 13--15 partial verdict/%
certificate marks are all model-theoretic. The two capabilities this work
combines---measured grounding and a feasibility verdict with a certificate---are held by
\emph{disjoint} sets of prior methods, and their intersection is empty. We compose the
measured genome-scale CRISPRi effect vectors from a Perturb-seq screen in primary human
CD4$^+$ T cells \citep{Zhu2025} directly, with no inferred network. Because knockdown is
loss-of-function and combinations are non-negative, reachability becomes membership in a
convex cone, solved by non-negative least squares \citep{LawsonHanson1974}: ``reachable''
returns the mixing weights (a minimal ranked recipe), and ``outside the cone'' is a theorem
whose separating hyperplane---a Farkas/KKT certificate \citep{Farkas1902}---names the genes
that would have to be \emph{activated} instead, a falsifiable CRISPRa hypothesis rather than
a null result. We validate by held-out-gene generalisation against a shuffled-target null
(Th2$\rightarrow$Th1 held-out cosine $0.448$, null $z\approx24$; KKT residual
$\approx\!10^{-11}$), and the operator transfers unchanged to the Norman K562 combinatorial
CRISPR screen \citep{Norman2019} (held-out CEBPA cosine $0.878$), whose measured
double-perturbations show that the additivity assumption is a bounded approximation (median
cosine $0.71$, sum-of-singles versus measured) rather than exact. We claim novelty for this
measured-effect, certificate-carrying reachability \emph{decision}, not for the T-cell or
CEBPA regulator lists, which the source screens \citep{Zhu2025,Norman2019} already report
and which we use only as positive controls; activation hypotheses require a separate CRISPRa
arm \citep{Schmidt2022}, and the modest absolute cosine is the ceiling of
loss-of-function alone.
